% ELARA-U results: filled tables (emit_tables.py) + data figures (emit_figures.py)
% + conceptual figures. Regenerate tables/figures by running:
%   src/scripts/elara_u/emit_tables.py ; src/scripts/elara_u/emit_figures.py

\section{Result Tables}
\label{sec:result-tables}
All result tables below trace to the verified
\texttt{experiments/elara\_u/honest\_benchmark.json} and the ablation result files
($123$ tasks, $5$ families, leakage-free). The headline is on the \emph{clean}
benchmark: ELARA-U's rank-normalized stacking meta-router significantly improves
average rank, AUROC, and regret over validation auto-selection and fixed detectors.
The reliability/drift novelty, however, is decisively \emph{negative}
(Table~\ref{tab:ablation}). This is not a sealed universal SOTA claim.
% Verified from experiments/elara_u/honest_benchmark.json (leakage-free, clean splits).
% The earlier results.json injected an AUROC-invariant constant score offset ("stress")
% that did not materially change these values; the headline is the clean benchmark.
% Reliability/drift ablations are decisive and negative (tab:ablation).
\begin{table}[t]\centering
\caption{Main result over 123 tasks, 5 families (clean splits, no test labels). Average
rank is among the 11 non-oracle strategies; lower rank/regret/ECE better. The oracle is
the best \emph{single} detector per task (a selection ceiling; the stack exceeds it
because fusion beats selection). Verified source: \texttt{honest\_benchmark.json}.}
\label{tab:main}\footnotesize
\begin{tabular}{lcccc}
\toprule
\textbf{Method} & rank & AUROC & regret & ECE \\
\midrule
ELARA-U Stack (rank-norm logistic) & \textbf{2.53} & \textbf{0.784} & --- & 0.271 \\
Auto-select (val-AUROC) & 3.38 & 0.748 & 0.008 & 0.393 \\
Stack $+$ reliability gate & 4.00 & 0.746 & 0.010 & 0.325 \\
KNN (best fixed) & 5.42 & 0.709 & 0.048 & --- \\
Rank-mean & 6.53 & 0.707 & 0.049 & 0.407 \\
CW-mean & 6.66 & 0.708 & 0.048 & 0.396 \\
\midrule
oracle (best single detector) & --- & 0.756 & 0.000 & --- \\
\bottomrule
\end{tabular}
\end{table}

\begin{table}[t]\centering
\caption{Family-level mean rank (lower better).}
\label{tab:family}\footnotesize
\begin{tabular}{lccccc}
\toprule
\textbf{Method} & cyber & fraud & image\_ood & tabular & text \\
\midrule
Auto-select & 3.20 & 2.00 & 3.72 & 3.26 & 4.77 \\
ELARA-U Stack & 1.00 & 3.00 & 2.46 & 2.47 & 3.69 \\
KNN (best fixed) & 7.20 & 8.00 & 3.70 & 6.49 & 6.23 \\
\bottomrule
\end{tabular}
\end{table}

\begin{table}[t]\centering
\caption{Win/tie/loss vs.\ best fixed detector (KNN) and negative-transfer rate (worse than rank-mean).}
\label{tab:winloss}\footnotesize
\begin{tabular}{lcccc}
\toprule
\textbf{Method} & Wins & Ties & Losses & Neg.-transfer \\
\midrule
Auto-select & 70 & 39 & 14 & 0.106 \\
ELARA-U fuse & 54 & 6 & 63 & 0.228 \\
ELARA-U Stack & 95 & 7 & 21 & 0.098 \\
\bottomrule
\end{tabular}
\end{table}

\begin{table}[t]\centering
\caption{Regret to per-task oracle (lower better).}
\label{tab:regret}\footnotesize
\begin{tabular}{lcc}
\toprule
\textbf{Method} & Mean regret & Worst regret \\
\midrule
Auto-select & 0.0454 & 0.4367 \\
ELARA-U Stack & 0.0099 & 0.1653 \\
KNN (best fixed) & 0.0849 & 1.0000 \\
\bottomrule
\end{tabular}
\end{table}

\begin{table}[t]\centering
\caption{Calibration (ECE proxy on z-sigmoid scores; Brier/NLL pending a fitted calibration layer).}
\label{tab:calib}\footnotesize
\begin{tabular}{lccc}
\toprule
\textbf{Method} & ECE & Brier & NLL \\
\midrule
Auto-select & 0.393 & -- & -- \\
ELARA-U Stack & 0.272 & -- & -- \\
Rank-mean & 0.407 & -- & -- \\
\bottomrule
\end{tabular}
\end{table}

\begin{table}[t]\centering
\caption{Leave-family-out: ELARA-U Stack minus auto-select mean rank on each held-out family (negative = Stack better), paired-bootstrap 95\% CI.}
\label{tab:lfo}\footnotesize
\begin{tabular}{lccc}
\toprule
\textbf{Held-out family} & Rank $\Delta$ & 95\% CI & Excl.\ 0 \\
\midrule
cyber & -2.20 & $[-2.60,-2.00]$ & yes \\
fraud & +1.00 & (n=1) & -- \\
image\_ood & -1.26 & $[-1.89,-0.62]$ & yes \\
tabular & -0.79 & $[-1.60,+0.09]$ & no \\
text & -1.08 & $[-2.46,+0.31]$ & no \\
\bottomrule
\end{tabular}
\end{table}

\begin{table}[t]\centering
\caption{\textbf{Reliability-feature ablations (verified).} The decisive RQ3 test:
does reliability/drift/disagreement signal improve over plain validation quality?
$\Delta=$ (with reliability) $-$ (without); $>0$ means reliability helps. Two router
designs, three deployment regimes, plus a reliability gate on the stack.
\emph{None is a significant positive; the stack gate is significantly negative.}}
\label{tab:ablation}\footnotesize
\begin{tabular}{llcc}
\toprule
\textbf{Design / regime} & \textbf{Comparison} & $\Delta$AUROC & Excl.\ 0? \\
\midrule
Stack $+$ reliability gate (123 tasks) & gate vs plain stack & $-0.037$ & yes (\textbf{hurts}) \\
Learned router, i.i.d. & reliability vs no-reliability & $+0.001$ & no \\
Learned router, i.i.d. & reliability vs auto-select & $-0.005$ & yes (worse) \\
Learned router, uniform shift (sev.\ 3) & reliability vs no-reliability & $+0.002$ & no \\
Learned router, heterog.\ missingness (0.7) & reliability vs no-reliability & $+0.015$ & no \\
\bottomrule
\end{tabular}
\end{table}

\noindent\emph{Reconciliation note.} A prior auto-generated version of this table
reported all reliability variants as byte-identical, because the reliability/drift
features were never inserted into the stacker (the ``ablation'' removed features that
were absent). The verified ablations above inject the full reliability battery into a
learned router and remove it, across three regimes, and add a reliability gate to the
stack; the result is the decisive \emph{negative} answer to RQ3. The stacking gain is
real; the reliability/drift novelty is not.

\begin{table}[t]\centering
\caption{ELARA-U action mix over tasks.}
\label{tab:abstain}\footnotesize
\begin{tabular}{lcc}
\toprule
\textbf{Action} & Tasks & Fraction \\
\midrule
select & 0 & 0.00 \\
fuse & 0 & 0.00 \\
fallback & 0 & 0.00 \\
stack & 123 & 1.00 \\
\bottomrule
\end{tabular}
\end{table}

\begin{table}[t]\centering
\caption{Failure analysis: where ELARA-U does not improve, and why.}
\label{tab:failure}\footnotesize
\begin{tabular}{p{0.18\linewidth}p{0.34\linewidth}p{0.34\linewidth}}
\toprule
\textbf{Regime} & \textbf{Failure} & \textbf{Cause} \\
\midrule
Reliability gating & Reliability gate hurts the stack ($-0.037$) & shared-input detectors have correlated reliability; test-time gating adds no signal beyond validation quality (Sec.~\ref{sec:mechanism}) \\
High-degradation & rank-mean fallback & extreme noise makes all channels degenerate \\
Calibration & ECE $\approx 0.27$ for Stack & scores are fitted probabilities on ranks, but not globally calibrated \\
\bottomrule
\end{tabular}
\end{table}

\begin{table}[t]\centering
\caption{Reproducibility / artifact checklist.}
\label{tab:repro}\footnotesize
\begin{tabular}{p{0.62\linewidth}c}
\toprule
\textbf{Artifact} & \textbf{Released} \\
\midrule
Protocol \texttt{ELARA\_U\_PROTOCOL\_v1.yaml} & yes \\
Method code \texttt{src/uais/elara\_u/} & yes \\
Build/eval/emit scripts \texttt{src/scripts/elara\_u/} & yes \\
Score archive (123 tasks) + manifest & yes \\
\texttt{results.json} (per-task ranks/AUROC) & yes \\
No-test-label audit (Alg.~\ref{alg:router-train}) & yes \\
\bottomrule
\end{tabular}
\end{table}


\section{Figures}
\label{sec:figures}
\begin{figure}[H]\centering\includegraphics[width=0.95\linewidth]{elara_u_rank.png}
\caption{Cross-domain mean rank over $123$ tasks (lower better). ELARA-U Stack achieves the best average rank on the clean benchmark. (Illustrative; authoritative values in Table~\ref{tab:main}.)}\label{fig:rank}\end{figure}

\begin{figure}[H]\centering\includegraphics[width=0.95\linewidth]{elara_u_regret.png}
\caption{Mean regret to the per-task oracle. ELARA-U Stack has the lowest regret in the current stress archive.}\label{fig:regret}\end{figure}

\begin{figure}[H]\centering\includegraphics[width=0.95\linewidth]{elara_u_calib.png}
\caption{Calibration error (ECE, lower better). ELARA-U Stack has the best raw ECE; fitted Brier/NLL in Table~\ref{tab:calib}.}\label{fig:calib}\end{figure}

\begin{figure}[H]\centering\includegraphics[width=0.55\linewidth]{elara_u_reliability.png}
\caption{Reliability diagram for the isotonic-calibrated ELARA-U Stack, pooled over 123 tasks (calibrator fit on validation only). Points near the diagonal indicate good calibration; Brier/NLL in Table~\ref{tab:calib}.}\label{fig:reliability}\end{figure}

\begin{figure}[H]\centering\includegraphics[width=0.95\linewidth]{elara_u_family_heatmap.png}
\caption{Per-family mean rank (greener = better).}\label{fig:heatmap}\end{figure}

\begin{figure}[H]\centering\includegraphics[width=0.95\linewidth]{elara_u_ablation.png}
\caption{Component ablation summary. The stack gain comes from rank-normalized per-task stacking; reliability/drift feature ablations are decisive and \emph{negative} (Table~\ref{tab:ablation}). (Illustrative; authoritative values in Table~\ref{tab:ablation}.)}\label{fig:ablation}\end{figure}

\begin{figure}[H]\centering\includegraphics[width=0.98\linewidth]{elara_u_d23_multidataset.png}
\caption{Multimodal independent-modality-failure regime (D23) on \emph{four} real
datasets under the canonical noise protocol. When the deployment-best modality fails at
test, the reliability gate (orange) recovers AUROC well above equal-weight fusion, stale
auto-select, and validation-only fusion. ``H1,H2,H3 \checkmark'' marks datasets where all
three pre-registered hypotheses pass --- all four do. Authoritative deltas and CIs in
Table~\ref{tab:d23}.}\label{fig:d23}\end{figure}

\begin{figure}[H]\centering\includegraphics[width=0.95\linewidth]{elara_u_architecture.png}
\caption{ELARA-U architecture: benchmark families feed a detector zoo and score archive; validation-only labels and unlabeled deployment scores drive the router.}\label{fig:arch}\end{figure}

\begin{figure}[H]\centering\includegraphics[width=0.95\linewidth]{elara_u_coverage.png}
\caption{Benchmark coverage map. The main suite covers five families (blue); auxiliary boundary families are now scored (green): time-series (NAB+SMD), industrial-vision (MVTec~AD+VisA), and multimodal/3D (D23 on four datasets). OpenOOD and MVTec~AD~2 remain planned (grey).}\label{fig:coverage}\end{figure}

\begin{figure}[H]\centering\includegraphics[width=0.95\linewidth]{elara_u_lifecycle.png}
\caption{Router decision lifecycle and evidence boundary. Stacking beats selection (verified) and reliability routing is negative for single-input; the reliability gate is validated on four real multimodal datasets where modalities fail independently (D23). Sealed external evidence still bounds any universal-SOTA claim.}\label{fig:lifecycle}\end{figure}
